
Standard LoRA targets linear projections but cannot directly modify SSM core modules such as the A, B, C, and D matrices that govern state dynamics. Recent work has cataloged this limitation and proposed alternatives including Sparse Dimension Tuning and State-offset Tuning. However, the question of when projection-only LoRA will succeed versus fail on SSM architectures has lacked a principled, measurable criterion.

This work investigates the spectral dynamics underlying SSM state evolution. SSM state dynamics are governed by the discretized transition matrix A\_bar = exp(Delta * A), whose eigenvalues determine information decay rates. The spectral memory horizon, defined as H\_spec = -1/log|lambda\_max|, gives the theoretical persistence length of the slowest-decaying eigenmode. Projection-only LoRA modifies the input and output mappings but leaves A\_bar eigenvalues unchanged.

This structural separation motivates two competing hypotheses:

\begin{enumerate}
\item \textbf{Memory Horizon Separation Hypothesis (MHSH):} Task success is bounded by whether information dependencies fall within H\_spec. Beyond-horizon tasks require modifying the A matrix eigenvalues.
\end{enumerate}

\begin{enumerate}
\item \textbf{Eigenmode Utilization Hypothesis (EUH):} Projection-only LoRA can succeed on beyond-horizon tasks by redistributing state energy toward slow eigenmodes, utilizing latent memory capacity.
\end{enumerate}

The experiments in this work test these hypotheses through measurement of eigenvalue preservation and energy redistribution under LoRA training. The main findings are:

\begin{itemize}
\item H\_spec is a stable, input-independent property measurable from pretrained Mamba weights (CV = 2.22 x $10^-16$).
\item Projection-only LoRA achieves exact eigenvalue preservation (delta H\_spec = 0.0\%, eigenvalue correlation = 1.0).
\item Energy redistribution toward slow eigenmodes is negligible (delta E = 5.93 x $10^-7$ nats), eliminating the EUH mechanism.
\end{itemize}

These results provide evidence that projection-only LoRA is bounded by the intrinsic spectral horizon, though the direct task-failure verification (H-M4) was not completed.
